% Este fuente se publica en dos dependencias causales mediante los
% selectores definidos por los wrappers de integración. Ambos tramos cubren
% conjuntamente el archivo completo, sin solape ni omisión.
\ifdefined\HMTGMCOperatorialAPP
\section[Matrix range of APP]{Matrix range of APP: exact classification
of the sum--product algebra and free interval}
\label{p12-83-c22-s1:gmc:chap:operatorios}

The qutrit squares of \cref{p12-83-c22-s3:gmc:chap:qms} provide an
operatorial input. There exists a
second input, more internal to APP: the two partitions induced by sum and
product generate noncommutative orthogonal projections. The language of
operator systems makes it possible to study them without forcing a grid.

\subsection{Pair of conditional expectations associated with the additive and multiplicative partitions of APP}

Let
\[
 \mathcal H_d=\ell^2(\{1,\ldots,9\}^d),
\]
and let us define
\[
 \Sigma_d(x)=\operatorname{dr}(x_1+\cdots+x_d),
 \qquad
 \Pi_d(x)=\operatorname{dr}(x_1\cdots x_d).
 \tag{GMC-6.1}
\]
Let $P_{\Sigma,d}$ and $P_{\Pi,d}$ be the orthogonal projections onto the
functions constant on the fibres of $\Sigma_d$ and $\Pi_d$. The preceding
bridge proves
\[
 \|[P_{\Sigma,2},P_{\Pi,2}]\|=\frac{\sqrt2}{3},
 \qquad
 \|[P_{\Sigma,3},P_{\Pi,3}]\|=\frac{\sqrt3}{4},
 \tag{GMC-6.2}
\]
and
\[
 \operatorname{Ran}P_{\Sigma,d}\cap
 \operatorname{Ran}P_{\Pi,d}=\mathbb C\mathbf1
 \qquad(2\le d\le15).
 \tag{GMC-6.3}
\]
APP therefore contains a noncommutative operator package prior to
any QLS or QMS designation.

We define
\[
 \mathcal S_d=\Span\{I,P_{\Sigma,d},P_{\Pi,d}\},
 \qquad
 \mathfrak C_d^{\rm APP}=C^*(P_{\Sigma,d},P_{\Pi,d}).
 \tag{GMC-6.4}
\]
$\mathcal S_d$ is a system of operators and its rango matrix is
\[
 \mathcal W_s(\mathrm{APP}_d)=
 \left\{(\Phi(P_{\Sigma,d}),\Phi(P_{\Pi,d})):
 \Phi:\mathfrak C_d^{\rm APP}\to M_s\ \text{UCP}\right\}.
 \tag{GMC-6.5}
\]
The tower $\bigsqcup_s\mathcal W_s$ is matricially convex. Unlike
a general label of “quantumness,” \textup{(GMC-6.5)} makes it possible to compute extreme points,
compressions, dilations, and invariants.

\subsection{Classification of the sum--product algebra at depth 2}

Place $P=P_{\Sigma,2}$ and $Q=P_{\Pi,2}$ on $\mathcal H_2$, of dimension
81. The nine fibres of $\Sigma_2$ have size nine. The fibres of
$\Pi_2$ have sizes
\[
 6,6,12,6,6,12,6,6,21.
 \tag{GMC-6.6}
\]
Let $u_a$ and $v_b$ be the normalized indicator vectors of those fibres, and
$C=(\langle u_a,v_b\rangle)_{a,b=1}^9$ their cross-Gram matrix. The exact
intersection count yields
\[
 \chi_{CC^*}(t)
 =t^5(t-1)\left(t-\frac57\right)
       \left(t-\frac13\right)^2.
 \tag{GMC-6.7}
\]
Moreover,
\[
\begin{array}{c|c}
\text{common subspace} & \text{dimension}\\
\hline
\ker P\cap\ker Q&64\\
\operatorname{Ran}P\cap\operatorname{Ran}Q&1\\
\operatorname{Ran}P\cap\ker Q&5\\
\ker P\cap\operatorname{Ran}Q&5
\end{array}
\tag{GMC-6.8}
\]

\begin{theorem}[Sum--product algebra of APP in $d=2$]
\label{p12-83-c22-s1:gmc:thm:algebra-app-d2}
There is an isomorphism of finite-dimensional $C^*$-algebras
\[
 \boxed{\ 
 \mathfrak C_2^{\rm APP}
 \cong\mathbb C^4\oplus M_2(\mathbb C)\oplus M_2(\mathbb C).
 \ }
 \tag{GMC-6.9}
\]
In particular, $\dim\mathfrak C_2^{\rm APP}=12$.
\end{theorem}

\begin{proof}
The canonical Halmos decomposition for two projections separates the four
subspaces of \textup{(GMC-6.8)} and a generic sector for each principal angle
$0<\lambda<1$. Equation \textup{(GMC-6.7)} shows that the only distinct
generic values are $5/7$ and $1/3$. In each case, the pair is unitarily
equivalent to
\[
 P_0=\begin{pmatrix}1&0\\0&0\end{pmatrix},
 \qquad
 Q_\lambda=
 \begin{pmatrix}
  \lambda&\sqrt{\lambda(1-\lambda)}\\
  \sqrt{\lambda(1-\lambda)}&1-\lambda
 \end{pmatrix},
 \tag{GMC-6.10}
\]
which generates $M_2$. The multiplicity of $\lambda=1/3$ amplifies the same
representation but does not create another simple summand. The four sectors
$0/1$ simultaneously produce $\mathbb C^4$.
\end{proof}

Specifically,
\[
 Q_{5/7}=
 \begin{pmatrix}5/7&\sqrt{10}/7\\\sqrt{10}/7&2/7\end{pmatrix},
 \qquad
 Q_{1/3}=
 \begin{pmatrix}1/3&\sqrt2/3\\\sqrt2/3&2/3\end{pmatrix}.
 \tag{GMC-6.11}
\]
The norm of the commutator is the maximum of
$\sqrt{\lambda(1-\lambda)}$ and recovers $\sqrt2/3$.

\begin{corollary}[Complete parameterization of the matrix range]
\label{p12-83-c22-s1:gmc:cor:rango-app-d2}
A pair $(X,Y)$ belongs to $\mathcal W_s(\mathrm{APP}_2)$ if and only if
there exist effects
$E_{00},E_{11},E_{10},E_{01}\succeq0$ and CP maps
$\Psi_{5/7},\Psi_{1/3}:M_2\to M_s$ such that
\[
 E_{00}+E_{11}+E_{10}+E_{01}
 +\Psi_{5/7}(I)+\Psi_{1/3}(I)=I_s,
 \tag{GMC-6.12}
\]
\[
 X=E_{11}+E_{10}
 +\Psi_{5/7}(P_0)+\Psi_{1/3}(P_0),
 \tag{GMC-6.13}
\]
\[
 Y=E_{11}+E_{01}
 +\Psi_{5/7}(Q_{5/7})+\Psi_{1/3}(Q_{1/3}).
 \tag{GMC-6.14}
\]
\end{corollary}

\begin{proof}
Every CP map from a direct sum decomposes into CP maps on its
summands. The four scalar summands produce the effects $E_{ab}$; the
two matrix summands produce $\Psi_\lambda$. Unitality is
\textup{(GMC-6.12)}; the coordinates of $P$ and $Q$ in \textup{(GMC-6.9)} give
\textup{(GMC-6.13)}--\textup{(GMC-6.14)}. The argument is reversible.
\end{proof}

\begin{bridgebox}
The classification \textup{(GMC-6.9)} provides the operator-theoretic realization of the bridge. It converts the
sum--product incompatibility into a fully computable operator object:
four classical sectors and two irreducible matrix blocks.
The value $1/3$ reappears here as a principal operator angle, not merely as
an MUB overlap; $5/7$ is the second generic value of $PQP$ in the
sum--product sector. Neither should be identified with the contraction $5/9$ of the
central block $B_c$. They are invariants of distinct objects, coordinated without
being identified by decimal similarity.
\end{bridgebox}

\subsection{Spectral separation \texorpdfstring{$4$}{4} and free matrix interval}

For the block of \textup{(GMC-2.4)},
$B_c=7I+2R$ with $R=R^*$ nonscalar, $R^2=I$ and
$\operatorname{Spec}(R)=\{-1,+1\}$,
$C^*(B_c)=C^*(R)\cong\mathbb C^2$. Its matrix rank at level $s$ is
\[
 \mathcal W_s(B_c)=
 \{\Phi(B_c):\Phi:C^*(B_c)\to M_s\text{ UCP}\}.
\]

\begin{theorem}[Free interval of the central block]
\label{p12-83-c22-s1:gmc:thm:intervalo-bc}
For every $s\ge1$,
\[
 \boxed{\ 
 \mathcal W_s(B_c)
 =\{Y=Y^*:5I_s\preceq Y\preceq9I_s\}.
 \ }
 \tag{GMC-6.15}
\]
\end{theorem}

\begin{proof}
A UCP map sends the symmetry $R$ to a self-adjoint contraction
$-I_s\preceq X\preceq I_s$ and, therefore,
$\Phi(B_c)=7I_s+2X$ belongs to the interval. Conversely, given such an $X$,
the effects $(I_s\pm X)/2$ define a UCP map from $\mathbb C^2$ that sends
$R$ to $X$.
\end{proof}

Thus, $7$ is the center, $2$ the operatorial radius, and $4$ the spectral diameter
of the free interval. The positive boundary, the boundary of the matrix interval, and the
determinant-one hyperbola of \textup{(GMC-2.6)} are three distinct notions that
can now be coordinated without being conflated.

The double projection does not reside in this early classification. Its formulation
from a common spectral antecedent is published, with its hypotheses and limits,
in the dedicated section of the governing chapter on double projection.
\fi

\ifdefined\HMTGMCDoubleProjection
\subsection{Common antecedent of readers: spectral refinement and joint dilation limit}

Let $X_m$ be a finite set of histories and let
\[
 R_m:X_m\to\mathcal A_m,
 \qquad
 Q_m:X_m\to\mathcal E_m
 \tag{GMC-6.16}
\]
two readers: for example, one Archimedean and the other exceptional. Each one
defines a \PVM in $C(X_m)$ by means of the characteristic functions of its
fibers.

\begin{proposition}[Common spectral predecessor]
\label{p12-83-c22-s1:gmc:prop:doble-proyeccion-conmutativa}
The two \PVM of \textup{(GMC-6.16)} commute and possess the common refinement
\[
 P_{a,e}=1_{\{R_m=a\}}1_{\{Q_m=e\}}
 =1_{\{x:R_m(x)=a,\ Q_m(x)=e\}}.
 \tag{GMC-6.17}
\]
If the readers are compatible with the truncations, the refinements are
compatible in $m$ and determine two readers on $C(X_\infty)$ with a
common cylindrical \PVM.
\end{proposition}

\begin{proof}
All functions belong to the same commutative algebra. The product
of the indicators is the indicator of the intersection. Summing in $e$ recovers
the \PVM of $R_m$, and summing in $a$ recovers that of $Q_m$. The compatibility passes
to the limit by the \Cref{p12-83-c22-s4:gmc:thm:dilatacion-diagonal}.
\end{proof}

This result formalizes the common antecedent of the double projection. It does
not prove that the pair of readers is jointly injective: for that purpose, it must
separate points of $X_\infty$. Nor does it automatically construct a common
noncommutative antecedent.

\begin{proposition}[Incompatibility of Two Mutually Unbiased Projective Measurements]
Two mutually unbiased qutrit \PVM do not admit a joint sharp measurement
in $M_3(\mathbb C)$.
\end{proposition}

\begin{proof}
Two finite \PVM are jointly measurable by a common \PVM if and only if
they commute. Projectors from distinct MUB contexts have overlap $1/3$ and
do not commute.
\end{proof}


The HMT double projection must therefore not be identified with two sharp MUB
contexts on the same qutrit. It can be realized on the common diagonal of
histories, through noisy effects, or in a larger noncommutative antecedent
with two UCP channels. The last possibility requires inclusion of matrix
ranges, Choi positivity and equivariance with respect to TPK dynamics.

\begin{workbox}[fontupper=\footnotesize,fonttitle=\bfseries\footnotesize,
  boxsep=1pt,left=4pt,right=4pt,top=0pt,bottom=0pt,
  toptitle=1pt,bottomtitle=1pt]
The noncommutative antecedent is constructed in
\Cref{sec:cpe-antecedente-ucp}: the system $\mathcal S_\infty$ and the UCP maps
$\Phi_{\rm arq}$ and $\Phi_{\rm exc}$ simultaneously reproduce the two
HMT projections and intertwine the TPK dynamics on the declared equivariant
domain. The \Cref{p12-83-c22-s1:gmc:prop:doble-proyeccion-conmutativa}
provides its diagonal antecedent.
\end{workbox}
\fi


\section{Genealogical matrix convexity: objects, domains, and typed dictionary}

\subsection{Substructures, global object, and projections}

The question organizing the work of De las Cuevas--Netzer is how a
collection of local measurements can compose a two-dimensional object and to
what extent that object is the compression of an ideal totality. An isolated
measurement ---a \POVM--- admits a Naimark or Stinespring dilation.
However, a grid simultaneously contains one \POVM in every row and
every column. Requiring a common dilation that makes all those measurements
projective is a much stronger global condition. There are quantum magic
squares that do not satisfy it, even at the smallest relevant noncommutative
matrix level \citep{delascuevas2020qms}.

HMT formulates another problem of parts and wholes. The visible coordinate of a
state does not necessarily preserve the route that produced it. APP distinguishes the
sum and product evaluations on a common support; TRIT preserves orientation
and carry; TPK preserves path composition, return, holonomy, and memory.
The continuum arises by coherently prolonging finite histories, and the
double projection publishes two different shadows of the same
enriched state. Both theories therefore ask what is lost under compression.
The difference is that one organizes loss by matrix level and the other by
genealogical depth.

\begin{intuitionbox}
De las Cuevas asks: «Can I view this imperfect grid as the
compression of a larger projective grid?». HMT asks: «What complete
history was compressed to produce this output?». The bridge answers that
both questions can be formulated on the same family, with two independent
indices: depth $m$ and observation dimension $s$.
\end{intuitionbox}

\subsection{Quantum Squares, Operator Systems, and Matrix Convexity}

Let $n,s\ge1$. A quantum magic square of exterior size $n$ and
interior size $s$ is a matrix
\[
 A=(A_{ij})_{i,j=0}^{n-1}\in\Mat_n(\Her_s(\mathbb C))
\]
such that
\[
 A_{ij}\succeq0,\qquad
 \sum_{j=0}^{n-1}A_{ij}=I_s,\qquad
 \sum_{i=0}^{n-1}A_{ij}=I_s.
 \tag{GMC-1.1}
\]
Each row and each column is a \POVM. If all entries are projections,
$A$ is a quantum permutation matrix; if, moreover, those projections
commute, $A$ is a commutative quantum permutation matrix
\citep{delascuevas2020qms,delascuevas2023latin}.

To \QMS is \emph{semiclasico} if there exists to \POVM
\(\{Q_\sigma\}_{\sigma\in S_n}\) such that
\[
 A=\sum_{\sigma\in S_n}P_\sigma\otimes Q_\sigma,
 \tag{GMC-1.2}
\]
where $P_\sigma$ is the classical permutation matrix associated with
$\sigma$. The formula does not require the effects $Q_\sigma$ to commute. It states
that the entire grid is obtained by reordering a single primary measurement.

A quantum Latin square of order $n$ is a grid of unit
vectors $v_{ij}\in\mathbb C^n$ whose rows and columns are
orthonormal bases. The rank-one projectors
\[
 P_{ij}=|v_{ij}\rangle\langle v_{ij}|
\]
form a projective \QMS. The easy-class \QLS are those produced by a classical Latin
square and a fixed orthonormal basis. De las Cuevas, Netzer, and
Valentiner-Branth prove that the semiclassical \QLS are exactly those of
the easy class and that the matrix convex hull of the \QLS contains more than the
semiclassical squares \citep{delascuevas2023latin}.

The elementary convex matrix operation is
\[
 X=\sum_{r=1}^{k}V_r^*X^{(r)}V_r,
 \qquad
 \sum_{r=1}^{k}V_r^*V_r=I_s.
 \tag{GMC-1.3}
\]
It permits objects of different internal sizes to be combined. In operator
language, compressions and direct sums are organized by unital completely
positive maps. The set of all matrix values of an operator system is its
matrix range.

\begin{importbox}
The governing results used are: a \QMS dilates to a quantum
permutation matrix with commuting entries if and only if it is semiclassical;
for $n\ge3$ and $s\ge2$ there exist nonsemiclassical \QMS, and for every $n\ge3$ there exist
\QMS on $M_2(\mathbb C)$ that do not dilate to any quantum
permutation matrix \citep{delascuevas2020qms}. Simultaneous dilatability
is therefore a substantive property, not an automatic consequence of
each row and column dilating separately.
\end{importbox}

\subsection{Typed correspondence between the HMT–MD genealogical state and operator systems: depth, truncation, monodromy, and reading}

In HMT, a visible output does not exhaust the state. A finite section
carries, according to the domain, data of the form
\[
 x_m=(\text{prefix},\text{sheet},\text{orientation},\text{carry},
       \text{boundary},\text{route},\text{memory},\text{gauge}).
 \tag{GMC-1.4}
\]
The depth $m$ is not a Hilbert-space dimension. It counts how much of the
genealogy has been resolved. Two states may share a visible coordinate and
differ in $m$, in leaf, or in memory.

APP provides the nonadic orthogram and the sum and product sheets. TRIT
records the ternary orientation and the local carry. TPK is the topological
phase kernel and the category of routes that composes those data. In a
finite realization, the routes act by operators on a Hilbert space,
but the representation does not replace the genealogy: it transports it.

The language of HMT also contains three operations that are absent from
the bare definition of a \QMS:

\begin{enumerate}
  \item \emph{truncation}: forgetting the continuation beyond depth
  $m$, preserving the prefix and aggregating its prolongations;
  \item \emph{monodromy}: returning to the same visible phase with nontrivial
  memory;
  \item \emph{reading}: publishing a mathematical or physical coordinate without
  identifying it with the complete section.
\end{enumerate}

\Needspace{.74\textheight}
\subsection{Direct and inverse correspondences between the formalisms}

\begin{longtable}{@{}L{.23\textwidth}L{.29\textwidth}L{.37\textwidth}@{}}
\toprule
Matrix language & Lenguaje HMT & Correspondencia tipada\\
\midrule
\endfirsthead
\toprule
Lenguaje matricial & Lenguaje HMT & Correspondencia tipada\\
\midrule
\endhead
Outer size $n$ & published grid & number of rows and columns of the
magic object; it is not depth nor internal dimension\\
Internal size $s$ & carta operatorial & dimension of the effects or of the
matrix compression\\
Profundidad $m$ & genealogy & new index that does not reduce to $n$ or
$s$\\
\POVM & family of positive readers & effects that sum to the unit; the sum
expresses normalization, not complete memory\\
\PVM & contexto espectral & family of orthogonal projections that resolves
the identity\\
\UCP & matrix reader & map that preserves unitarity and positivity at all levels\\
Compression $V^*(\cdot)V$ & shadow & step from a higher realization to an
observational chart lower\\
Dilation & antecedente operatorio & common multiplicative representation
before compressing\\
Semiclassicality & a single reordered primary \POVM & formal property; does not
equivalent to commutativity of all cells\\
No semiclasicidad & genuinely two-dimensional composition & cannot be
reduced to the reordering of a common primary \POVM\\
System of operators & space of readers & self-adjoint subspace with the
unit that organizes all the UCP cards\\
Rango matricial & set of shadows to all the $s$ & free family of
internal evaluations\\
Inductive limit & closure of operatorial refinements & conservation of the
unit through the entire nonadic filtration\\
Inverse limit & infinite section compatible & history that preserves all
its prefixes\\
\bottomrule
\end{longtable}

\subsection{Extension of the bigrading to 3 indices}

The magic-square formalism separates $(n,s)$. HMT adds $m$. The
total object must be written
\[
 \mathfrak M^{(n)}_{m,s},
 \tag{GMC-1.5}
\]
and not simply \(\mathfrak M_n\) or \(\mathfrak M_s\). Each index answers
a different question:

\begin{center}
\begin{tabularx}{.92\textwidth}{@{}c L{.25\textwidth}Y@{}}
\toprule
Index & Question & Associated operation\\
\midrule
$m$ & How much history is preserved? & refinement and truncation\\
$n$ & How is the grid composed? & rows, columns and permutations\\
$s$ & From which matrix level is it observed? & direct sum and compression UCP\\
\bottomrule
\end{tabularx}
\end{center}

For APP, $n=9$ is fixed. Then the family
$\{\mathfrak M^{(9)}_{m,s}\}_{m,s}$ is bigraded by $(m,s)$. If
APP, the four \QLS of order three, and the atlas of twelve directions are
compared, the index $n$ varies once more, and the correct structure is trigraded.

\begin{controlbox}
Calling it a “bigrading $(m,s)$” does not erase $n$: it means
that a family with fixed exterior size is being studied. Calling APP
“$9\times9$” likewise does not fix the interior size of its operators.
The internal matrices $6\times6$, the qutrit chart $s=3$, and the
exterior grid $n=9$ are objects of distinct types.
\end{controlbox}

\subsection{Bidirectionality of Mathematical Exchange}

The De las Cuevas--Netzer $\longrightarrow$ HMT direction contributes criteria that HMT had not
formulated with this degree of separation: cellwise positivity,
normalization of all rows and columns, semiclassicality, the boundary of the
matrix convex hull, simultaneous dilation, and Arveson extreme points. The
result is a set of tests and falsifiers for any claim of a
«quantum square».

The HMT $\longrightarrow$ De las Cuevas--Netzer direction adds structure to the object
being compressed: depth, route, nonsplitting of the carry, monodromy, dynamics,
gauge, double projection, and metrological readers. The problem ceases to ask
only whether a dilation exists and instead asks whether a
\emph{genealogically faithful} dilation exists: one that does not identify two distinct routes merely
because they produce the same visible effect.

\begin{bridgebox}
The new notion organizing the article is the \emph{simultaneous
genealogical dilation}: a common multiplicative representation that dilates not only several
rows and columns at fixed depth, but all compatible depths,
and that specifies which part of memory remains in the representation and which
part is lost in compression. The exact theorem is proved in
\Cref{p12-83-c22-s4:gmc:chap:bigraduacion}.
\end{bridgebox}


\section[Operatorial incidence systems]{Operator systems of
incidence: quantum Latin squares, qutrit contexts, and common
genealogical dilation}
\label{p12-83-c22-s3:gmc:chap:qms}

The outer APP object has order nine, but the number nine does not suffice
to make it a quantum magic square. One must exhibit 81 positive effects
and verify 18 operator normalizations. This chapter performs
that task for two shadows of APP and for the complete atlas of twelve
directions. It also determines its exact class in the De las
Cuevas--Netzer taxonomy.

\subsection{4 cyclic quantum Latin squares of easy operator class}

Let
\[
 \mathcal P_a=\{P_{a,0},P_{a,1},P_{a,2}\},
 \qquad a\in\mathbb P^1(\mathbb F_3),
 \tag{GMC-4.1}
\]
the four qutrit \PVM of the Weyl system. For each $a$, choose
unit vectors $v_{a,k}$ with
$P_{a,k}=|v_{a,k}\rangle\langle v_{a,k}|$ and define
\[
 L^{(a)}_{ij}=v_{a,i+j\bmod3}.
 \tag{GMC-4.2}
\]
Each row and each column traverses the basis
$(v_{a,0},v_{a,1},v_{a,2})$ once. Therefore, $L^{(a)}$ is a \QLS of order three.
Since it is obtained from the cyclic Latin table and from a fixed orthonormal basis, it is of easy class and its
projective \QMS is semiclassical
\citep{delascuevas2023latin}.

\begin{proposition}[MUB Obstruction to Latin Stacking]
\label{p12-83-c22-s3:gmc:prop:no-go-mub}
Two rows from distinct contexts
$\mathcal P_a$ and $\mathcal P_b$, $a\ne b$, cannot be taken as rows of the same \QLS of
order three while preserving their aligned columns.
\end{proposition}

\begin{proof}
The column condition of a \QLS would require orthogonality between the two
vectors of a column. But mutual unbiasedness gives
\[
 |\langle v_{a,k},v_{b,\ell}\rangle|^2
 =\Tr(P_{a,k}P_{b,\ell})=\frac13
\]
for every $k,\ell$. The overlap is not zero.
\end{proof}

The obstruction does not deny quantum compatibility of the four bases.
It determines that the operator structure appropriate for using them simultaneously is not a
\QLS of rank one: one must pass to weighted effects and to a \QMS.

\subsection{Cyclic lemma for quantum magic squares associated with a positive operator-valued measure}

\begin{lemma}[Cyclic square associated with a \POVM]
\label{p12-83-c22-s3:gmc:lem:qms-ciclico}
Let $E=(E_0,\ldots,E_{n-1})$ be a \POVM on $\mathbb C^s$. Then
\[
 A_{ij}=E_{i+j\bmod n},
 \qquad 0\le i,j<n,
 \tag{GMC-4.3}
\]
is a \QMS of sizes $(n,s)$. It is also semiclassical and admits a
simultaneous dilation to a commutative \QPM.
\end{lemma}

\begin{proof}
Each entry is positive. Every row and every column is a permutation of the
$n$ effects, so that its sum is $I_s$. For $t\in\mathbb Z/n\mathbb Z$,
let $\sigma_t(i)=t-i$ and $P_{\sigma_t}$ be its permutation matrix. Then
\[
 A=\sum_{t=0}^{n-1}P_{\sigma_t}\otimes E_t,
 \tag{GMC-4.4}
\]
which is an explicit semiclassical decomposition. If
$E_t=V^*Q_tV$ is a Naimark dilation of the \POVM to a \PVM
$(Q_t)_t$, the entries
$\widehat A_{ij}=Q_{i+j}$ form a commutative \QPM and
$A_{ij}=V^*\widehat A_{ij}V$ simultaneously for the $n^2$ cells.
\end{proof}

This construction belongs to the class of Latin families associated with POVMs
(\emph{POVM Latin-square friends})
of De las Cuevas and Netzer \citep{delascuevas2023latin}. The novelty of the bridge lies not in the
abstract lemma, but in the \POVM furnished by HMT, their nonadic calibration, and the
exact separation between what each square preserves and what it still
compresses.

\subsection{Additive APP representation of sizes \((9,3)\) and its incidence codomain}

Fix one of the four contexts, for example
$\mathcal P_0=\{P_{0,0},P_{0,1},P_{0,2}\}$. For $t\in\mathbb Z/9\mathbb Z$
define
\[
 E_t^{\Sigma}=\frac13P_{0,t\bmod3},
 \qquad
 A_{ij}^{\Sigma}=E_{i+j\bmod9}^{\Sigma}.
 \tag{GMC-4.5}
\]
Each projection appears three times, then
\[
 \sum_{t=0}^{8}E_t^{\Sigma}
 =\frac13\sum_{k=0}^2 3P_{0,k}=I_3.
\]
By \Cref{p12-83-c22-s3:gmc:lem:qms-ciclico}, $A^\Sigma$ is a semiclassical \QMS of sizes
$(9,3)$. It is a direct transduction of the Latin law
$\Sigma(i,j)=i+j\bmod9$ into a qutrit reading. It preserves the outer size and
the additive leaf; it deliberately does not distinguish the three Hensel lifts of
the same ternary class.

\subsection{\(3\) contexts in a square $(9,3)$}

Elijamos three of the four \PVM qutrit, indexadas by $a=0,1,2$, and ordenemos
their projectors by means of $t=3a+k$:
\[
 E_{3a+k}^{(9)}=\frac13P_{a,k},
 \qquad a,k\in\mathbb Z/3\mathbb Z.
 \tag{GMC-4.6}
\]
Since
\[
 \sum_{a,k}E_{3a+k}^{(9)}
 =\frac13\sum_{a=0}^2\sum_{k=0}^2P_{a,k}=I_3,
\]
the formula
\[
 A_{ij}^{(9)}=E_{i+j\bmod9}^{(9)}
 \tag{GMC-4.7}
\]
produce otro \QMS of sizes $(9,3)$.

\begin{theorem}[Qutrit square of outer order nine]
\label{p12-83-c22-s3:gmc:thm:qms-9-3}
The square $A^{(9)}$ of \textup{(GMC-4.7)} satisfies:

\begin{enumerate}
 \item its 81 cells are positive effects of rank one and trace $1/3$;
 \item all rows and columns sum to $I_3$;
 \item contains entries that do not commute;
 \item is semiclassical by means of
 \[
 A^{(9)}=\sum_{t=0}^{8}P_{\sigma_t}\otimes E_t^{(9)};
 \tag{GMC-4.8}
 \]
 \item admits a simultaneous commutative dilation of its 81
 entries.
\end{enumerate}
\end{theorem}

\begin{proof}
The first two points and the decomposition follow from
\Cref{p12-83-c22-s3:gmc:lem:qms-ciclico}. If $a\ne b$, two rank-one MUB projectors do not
commute unless they coincide or are orthogonal, both of which are impossible because their overlap
is $1/3$; hence some
$[E_{3a+k}^{(9)},E_{3b+\ell}^{(9)}]$ are nonzero. The last assertion is the
dilation constructed in the proof of the lemma.
\end{proof}

\begin{bridgebox}
This theorem closes an existence that the preliminary bridge had left
open: there already are 81 explicit effects with $n=9$, $s=3$ and exact
normalizations. At the same time it proves a fertile conceptual separation:
\[
 \text{noncommutative cells}
 \quad\centernot\Longrightarrow\quad
 \text{no semiclassicality}.
\]
Here the cells do not commute, but the entire grid arises by reordering a
single primary \POVM.
\end{bridgebox}

\subsection{The atlas complete in a cuadrado \texorpdfstring{$(12,3)$}{(12,3)}}

The projective line over the nonadic ring has twelve points:
\[
 |\mathbb P^1(\mathbb Z/9\mathbb Z)|=12.
 \tag{GMC-4.9}
\]
Reduction modulo three induces
\[
 r:\mathbb P^1(\mathbb Z/9\mathbb Z)
 \longrightarrow\mathbb P^1(\mathbb F_3),
 \tag{GMC-4.10}
\]
with four fibers of cardinality three. Fixing a calibration
\[
 \Xi:\mathbb P^1(\mathbb Z/9\mathbb Z)
 \xrightarrow{\sim}
 \{(a,k):a\in\mathbb P^1(\mathbb F_3),\ k\in\mathbb F_3\}
 \tag{GMC-4.11}
\]
that respects the fibre of $r$, each direction $\ell$ is assigned the projector
$P_{\Xi(\ell)}$. Ordering the twelve directions as $t=0,\ldots,11$, let
\[
 F_t=\frac14P_{\Xi(\ell_t)},
 \qquad
 B_{ij}=F_{i+j\bmod12}.
 \tag{GMC-4.12}
\]

\begin{theorem}[Square qutrit of the complete nonadic atlas]
\label{p12-83-c22-s3:gmc:thm:qms-12-3}
The matrix $B$ is a \QMS of tamanos $(12,3)$, contains the twelve projectors
calibrados of the four contextos qutrit and is semiclasica:
\[
 B=\sum_{t=0}^{11}P_{\sigma_t}\otimes F_t.
 \tag{GMC-4.13}
\]
It admits a simultaneous dilation to a commutative \QPM.
\end{theorem}

\begin{proof}
Each of the four contexts sums to $I_3$, so
$\sum_tF_t=(1/4)\cdot4I_3=I_3$. The
\Cref{p12-83-c22-s3:gmc:lem:qms-ciclico} applies.
\end{proof}

The factor $1/4$ is not decorative. The atlas contains four resolutions of the
identity and, without that normalization, the sum would be $4I_3$. The complete atlas
is therefore a twelve-outcome measurement, not a fifth orthonormal
basis.

\subsection{Operatorial theorems and genealogical fidelity criterion APP--TPK}

The three previous constructions have an exact status:

\begin{center}
\begin{tabularx}{.96\textwidth}{@{}L{.20\textwidth}L{.21\textwidth}Y@{}}
\toprule
Object & Closed theorem & HMT Structure Not Yet Transported\\
\midrule
$A^\Sigma$, $(9,3)$ & sheet sum, positivity, normalization and
semiclassicality & product sheet, quotient and paths\\
$A^{(9)}$, $(9,3)$ & three contexts, noncommutativity and common dilation &
fourth context, Hensel incidence, gnomon and TPK memory\\
$B$, $(12,3)$ & complete atlas $12\to4\times3$ and common dilation &
equivariance of calibration, dynamics, and monodromy\\
\bottomrule
\end{tabularx}
\end{center}

Neither of the final two squares is a \QLS or a \QPM: their cells are projectors scaled by $1/3$ or $1/4$, not projections. They are quantum magic squares under the orthodox definition, and they do possess a simultaneous dilation. This formulation resolves the terminological ambiguity and preserves the mathematical difference between an effect and a projection.

\begin{workbox}
The next objective is not to prove again that a \QMS of order
nine exists. It is to construct an \emph{APP--TPK-faithful} one: in addition to \textup{(GMC-4.3)}, it must
simultaneously carry $\Sigma$, $\Pi$, quotient, orientation, and route, and it must
satisfy an equivariance equation with the TPK action. A candidate
that does not satisfy that equation will remain a valid shadow, but not the total
object.
\end{workbox}


\section[Genealogical--Matrix Bigrading]{Genealogical--matrix
bigrading at fixed outer size and trigrading
\texorpdfstring{$(m,n,s)$}{(m,n,s)}: truncation, compression, and simultaneous
dilation}
\label{p12-83-c22-s4:gmc:chap:bigraduacion}

The De las Cuevas--Netzer separation between outer size $n$ and matrix
level $s$ resolves an ambiguity in the language of grids. HMT requires
a second separation: depth $m$ is not an inner size. This
chapter proves that the $m$ and $s$ directions are compatible and that a
family coherent at every depth possesses a common dilation.

\subsection{Inverse system of histories and direct system of observables}

Let $(X_m,p_m)_{m\ge0}$ be a system of nonempty finite sets with surjective maps
\[
 p_m:X_{m+1}\longrightarrow X_m.
 \tag{GMC-3.1}
\]
In HMT, $X_m$ may be the complete set of nonadic prefixes of
length $m$ or a set of survivors, provided that the restriction maps
are surjective on the declared domain. The inverse limit
\[
 X_\infty=\varprojlim_m X_m
\]
is the compact space of compatible histories. Each $x\in X_m$ determines
the clopen cylinder
\[
 [x]_m=\{\xi\in X_\infty:\xi_m=x\}.
 \tag{GMC-3.2}
\]

By duality, the diagonal algebras
\[
 D_m=C(X_m)
\]
form an inductive system through
\[
 \jmath_m:D_m\longrightarrow D_{m+1},
 \qquad \jmath_m(f)=f\circ p_m.
 \tag{GMC-3.3}
\]
Its uniform limit is $C(X_\infty)$: the cylinder functions form a
unital self-adjoint subalgebra that separates points and are dense by
Stone--Weierstrass.

\subsection{The family at \(2\) degrees}

For $s\ge1$, define the matrix state space
\[
 \mathcal K_{m,s}=\UCP(D_m,M_s(\mathbb C)).
 \tag{GMC-3.4}
\]
If $e_x\in D_m$ is the characteristic function of $x\in X_m$, every
$\Phi\in\mathcal K_{m,s}$ is equivalent to a \POVM
\[
 E_x^{(m)}=\Phi(e_x)\succeq0,
 \qquad
 \sum_{x\in X_m}E_x^{(m)}=I_s.
 \tag{GMC-3.5}
\]
The genealogical restriction is
\[
 \rho_m^s:\mathcal K_{m+1,s}\longrightarrow\mathcal K_{m,s},
 \qquad
 \rho_m^s(\Phi)=\Phi\circ\jmath_m.
 \tag{GMC-3.6}
\]
In terms of effects,
\[
 E_x^{(m)}=\sum_{y\in p_m^{-1}(x)}E_y^{(m+1)}.
 \tag{GMC-3.7}
\]
It is exactly the operation “forget the continuation but sum all its
prolongations.” It is neither an average nor an elimination of mass: it preserves the unit.

In the matrix direction, if $d\ge0$,
$\Phi_r\in\mathcal K_{d,s_r}$ and
$V_r:\mathbb C^s\to\mathbb C^{s_r}$ satisfy
$\sum_rV_r^*V_r=I_s$, their matrix convex combination is
\[
 \operatorname{mc}_{\boldsymbol V}(\Phi_1,\ldots,\Phi_k)(f)
 =\sum_{r=1}^kV_r^*\Phi_r(f)V_r
 \in\mathcal K_{d,s}.
 \tag{GMC-3.8}
\]

\begin{theorem}[Genealogy Switching and Matrix Convexity]
\label{p12-83-c22-s4:gmc:thm:conmutacion-ms}
For every family $\Phi_r\in\mathcal K_{m+1,s_r}$ and every compression as in
\textup{(GMC-3.8)},
\[
 \rho_m^s\!\left(
   \operatorname{mc}_{\boldsymbol V}(\Phi_1,\ldots,\Phi_k)
 \right)
 =
 \operatorname{mc}_{\boldsymbol V}\!\left(
   \rho_m^{s_1}(\Phi_1),\ldots,\rho_m^{s_k}(\Phi_k)
 \right).
 \tag{GMC-3.9}
\]
\end{theorem}

\begin{proof}
For $f\in D_m$, both sides are equal to
\[
 \sum_rV_r^*\Phi_r(\jmath_m f)V_r.
\]
The equality is functorial, not numerical: it holds simultaneously for all
matrix levels and for any system of truncations.
\end{proof}

\begin{bridgebox}
Equation \textup{(GMC-3.9)} is the core of the bigrading $(m,s)$. Refining the
history and then compressing produces exactly the same shadow as
compressing each reader and then forgetting the continuation. Depth and
matrix level are independent, but form a commutative square.
\end{bridgebox}

\begin{figure}[H]
\centering
\begin{tikzpicture}[node distance=2.5cm and 3.7cm,>=Latex]
\node[draw,rounded corners,fill=hmtlightblue,minimum width=3.3cm,
 minimum height=1cm] (a) {$\mathcal K_{m+1,s'}$};
\node[draw,rounded corners,fill=hmtlightblue,right=of a,minimum width=3.3cm,
 minimum height=1cm] (b) {$\mathcal K_{m+1,s}$};
\node[draw,rounded corners,fill=hmtlightgold,below=of a,minimum width=3.3cm,
 minimum height=1cm] (c) {$\mathcal K_{m,s'}$};
\node[draw,rounded corners,fill=hmtlightgold,right=of c,minimum width=3.3cm,
 minimum height=1cm] (d) {$\mathcal K_{m,s}$};
\draw[->,thick] (a)--node[above]{compression} (b);
\draw[->,thick] (a)--node[left]{truncation} (c);
\draw[->,thick] (b)--node[right]{truncation} (d);
\draw[->,thick] (c)--node[below]{compression} (d);
\node[hmtviolet] at ($(a)!0.5!(d)$) {commutes};
\end{tikzpicture}
\caption{The two independent directions of the bigraduation. The symbol
$s'$ may represent a direct sum of several internal levels.}
\label{p12-83-c22-s4:gmc:fig:cuadrado-ms}
\end{figure}

\subsection{Uniform dilation compatible with all depths of the projective system}

A family $(\Phi_m)_{m\ge0}$ is compatible if
\[
 \Phi_m=\Phi_{m+1}\circ\jmath_m.
 \tag{GMC-3.10}
\]
This is equivalent to requiring the aggregation law \textup{(GMC-3.7)} for all prefixes.

\begin{theorem}[Simultaneous genealogical dilation, diagonal corridor]
\label{p12-83-c22-s4:gmc:thm:dilatacion-diagonal}
Every compatible family
$\Phi_m\in\UCP(C(X_m),M_s)$ determines a unique map
\[
 \Phi_\infty:C(X_\infty)\longrightarrow M_s
 \tag{GMC-3.11}
\]
which coincides with $\Phi_m$ on each cylindrical algebra. There exist a Hilbert
space $\mathcal K$, an isometry $V:\mathbb C^s\to\mathcal K$, a
representation $\pi:C(X_\infty)\to B(\mathcal K)$, and its associated projective
spectral measure $P$ on $X_\infty$ such that
\[
 \Phi_\infty(f)=V^*\left(\int_{X_\infty}f(\xi)\,dP(\xi)\right)V.
 \tag{GMC-3.12}
\]
In particular, for every $m$ and $x\in X_m$,
\[
 E_x^{(m)}=V^*P([x]_m)V.
 \tag{GMC-3.13}
\]
A common \PVM simultaneously dilates all finite \POVM of the
genealogy. If the Stinespring dilation is chosen to be minimal, the triple
$(\pi,\mathcal K,V)$ is unique up to unitary equivalence; once
$\pi$ is fixed, its spectral measure $P$ is unique.
\end{theorem}

\begin{proof}
Compatibility defines a unital positive map on the union of the
cylindrical algebras. Because the domain is commutative, positivity implies
complete positivity, and the map is contractive. By density it
extends uniquely to $C(X_\infty)$. The Stinespring theorem yields
$\Phi_\infty(f)=V^*\pi(f)V$ and guarantees that its minimal realization is unique
up to unitary equivalence. Every representation of a commutative algebra
is the integral with respect to a unique \PVM $P$ associated with that representation.
Applying the formula to the characteristic
function of the cylinder $[x]_m$ gives \textup{(GMC-3.13)}.
\end{proof}

\begin{corollary}[Fidelity environmental realization]
The representation of \Cref{p12-83-c22-s4:gmc:thm:dilatacion-diagonal} may be chosen faithful on
$C(X_\infty)$ without changing any finite compression.
\end{corollary}

\begin{proof}
If $(\pi,\mathcal K,V)$ is a dilation and
$\pi_f:C(X_\infty)\to B(\mathcal K_f)$ a faithful representation, take
$\pi\oplus\pi_f$ and the isometry $V\oplus0$. The compression remains
$\Phi_\infty$, and the ambient representation no longer identifies two distinct
cylindrical functions.
\end{proof}

The final construction types the phrase “preserve the genealogy.” The minimal dilation preserves exactly what the reader sees; the faithful ambient dilation additionally preserves all distinctions of the algebra before compression. The fidelity of a TPK dynamics, however, also requires intertwining its arrows and does not follow solely from the fidelity of the diagonal representation.

\subsection{Noncommutative UHF Tower and Matrix Compatibility}

The nonadic UHF tower, its trace, and its factorial closures are proved in
\cref{chap:integracion-vn64-operatorial}. Here they are taken as an explicit
starting structure to prove a different result: the simultaneous prolongation
of a compatible family of UCP maps and their common
dilation.

Let
\[
 \mathcal A_m=M_{9^m}(\mathbb C),
 \qquad \iota_m(a)=a\otimes I_9,
 \qquad
 \mathcal A_\infty=\varinjlim(\mathcal A_m,\iota_m).
 \tag{GMC-3.14}
\]

\begin{theorem}[Simultaneous dilation on the UHF limit]
\label{p12-83-c22-s4:gmc:thm:dilatacion-uhf}
If $\Psi_m:\mathcal A_m\to M_s$ are UCP maps and
\[
 \Psi_m=\Psi_{m+1}\circ\iota_m,
 \tag{GMC-3.15}
\]
there exists a unique UCP map
$\Psi_\infty:\mathcal A_\infty\to M_s$ that prolongs them. There exists a
representation $\pi:\mathcal A_\infty\to B(\mathcal K)$ and an isometry
$V:\mathbb C^s\to\mathcal K$ such that
\[
 \Psi_m(a)=V^*\pi(\iota_{m,\infty}(a))V
 \quad\text{for every }m.
 \tag{GMC-3.16}
\]
The representation can be made faithful by direct sum with a faithful representation of the UHF limit.
\end{theorem}

\begin{proof}
Compatibility defines $\Psi_\infty$ on the algebraic inductive limit. UCP
maps are completely contractive, so it extends to the
completion $C^*$. Stinespring gives \textup{(GMC-3.16)}; the direct sum proves the
last assertion.
\end{proof}

Restricting \Cref{p12-83-c22-s4:gmc:thm:dilatacion-uhf} to the canonical diagonals recovers
the complete nonadic case $X_m=(\mathbb Z/9\mathbb Z)^m$ of the
\Cref{p12-83-c22-s4:gmc:thm:dilatacion-diagonal}; this restriction does not by itself cover an
arbitrary system of survivors. In the canonical nonadic UHF system of
\textup{(GMC-3.14)}, the GNS representation of its unique trace has weak closure
isomorphic to $\Rhyper$. Thus, the bigrading is not a combinatorial appendix: it joins
the cylinders of histories, the finite matrix charts, and the operatorial
continuum of the von Neumann realization.

\subsection{Trigraduation and Levels of Simultaneity}

When a grid appears, the exterior size $n$ is incorporated:
\[
 \mathfrak M_{m,s}^{(n)}.
 \tag{GMC-3.17}
\]
The structure is trigraded if $n$ varies and bigraded by $(m,s)$ if
$n$ is fixed. ‘Simultaneous’ can denote three increasingly stringent requirements:

\begin{enumerate}
 \item \emph{by rows and
columns}: a common dilation of the $2n$
\POVM of a magic square;
 \item \emph{by depth}: a common dilation of all levels
 $m$, proved in Theorems~\ref{p12-83-c22-s4:gmc:thm:dilatacion-diagonal}
 and~\ref{p12-83-c22-s4:gmc:thm:dilatacion-uhf};
 \item \emph{by dynamic genealogy}: in addition, an intertwiner preserving
 the route, monodromy, and TPK composition.
\end{enumerate}

The first two have exact theorems within their respective domains. The third is the
strongest notion that HMT contributes to the De las Cuevas--Netzer framework. The
construction of the common operator system and its equivariance equation are
proved in \Cref{sec:cpe-antecedente-ucp}; the static dilation constitutes
only one of its antecedents.

\begin{controlbox}
The theorem does not claim that every disconnected collection of effects at different
depths has a common dilation. The decisive hypothesis is the
compatibility \textup{(GMC-3.10)} or \textup{(GMC-3.15)}. In HMT, that compatibility is the
operator form of “each prefix is exactly the restriction of its
continuations”.
\end{controlbox}


\section{Matrix cones, ranks, and completely positive extensions}
\label{p12-83-c22-s5:gmc:chap:beyond}

The recent program of De les Coves, van der Eyden, and Netzer extends
operator systems to conic systems defined over more general
categories \citep{delescoves2026beyond}. The framework does not turn every
category into a quantum theory: it requires choosing suitable cones, dualities, compressions, and
completely positive maps. HMT provides an especially
rich testbed in which to make that choice.

\subsection{Realizations in systems of ordinary operators}

The following families are entirely within the standard formalism:

\begin{itemize}
 \item $\mathcal S_d=\Span\{I,P_{\Sigma,d},P_{\Pi,d}\}$ and its matrix
 ranks;
 \item the diagonal algebras $C(X_m)$ and the UHF tower
 $M_{9^m}\hookrightarrow M_{9^{m+1}}$;
 \item the readers UCP, \POVM, \PVM, channels CPTP and dilatacionis of
 Stinispring;
 \item the bigraded families $\mathcal K_{m,s}$.
\end{itemize}

No broader categorical formalism is required here in order to assert
positivity, dilation, or matrix convexity. Indeed, the classification
\textup{(GMC-6.9)} shows that a substantial part of APP can be resolved by
exact finite-dimensional operator theory.

\subsection{Genealogical extension of operator systems for TPK routes}

A TPK route is not merely an operator between two spaces: it carries source, target,
sheet, orientation, length, carry, and a composition law. If those
labels are forgotten and only the matrix representing the route is retained,
distinct arrows may collapse. To incorporate the complete type, one proposes
a category $\mathsf C_{\rm HMT}$ whose objects are register profiles and
whose arrows are admissible transports. To each type $c$ one would associate a
cone of positive readers $A(c)$, and to each transformation a
compatible action on those cones.

In the language of \emph{Beyond Operator Systems}, the objective is to construct a $S$-system: a functorial extension of the tensor product by the base object to a family of cones, with CP maps defined as natural transformations level by level. The theory provides minimal and maximal systems, duality, generalized Choi--Kraus theorems, representations as intersections of $S$-hedra, and a $*$-autonomous category \citep{delescoves2026beyond}.

\begin{workbox}
The intrinsic $S$-system of the TPK remains to be constructed. At a minimum, the
cone functor, the dualizing object, the evaluations and coevaluations, and the
proof that TPK composition preserves positivity in every type are lacking.
The fact that the TPK is monoidal in its finite realization does not by itself prove
that it is $*$-autonomous.
\end{workbox}

\subsection{Genealogical HMT Cones and Extension of the Conical Theory}

The classical examples start from a base cone and study its minimal and
maximal systems. HMT suggests a base cone with stratified memory:
\[
 \mathcal C_m^{\rm HMT}
 =\operatorname{cone}\{\text{positive readers of histories at depth }m\}.
 \tag{GMC-8.1}
\]
The truncations induce positive maps
$\mathcal C_{m+1}^{\rm HMT}\to\mathcal C_m^{\rm HMT}$, while the
compressions generate the matrix levels. The new phenomenon to study is
whether the minimal and maximal systems remain separate upon taking the limit in
$m$, and what part of that separation measures loss of genealogy rather than mere
noncommutativity.

This proposes three invariants for the general program:

\begin{enumerate}
 \item \emph{separation depth}: the smallest $m$ at which the minimal
 and maximal systems differ;
 \item \emph{matrix detection level}: the least $s$ exhibiting that
 difference;
 \item \emph{memory defect}: the minimal route information that must be
 added in order to lift a point from the maximal system to the minimal system.
\end{enumerate}

The first two refine the separation $(m,s)$; the third is genuinely
HMT. They are proposed definitions and must first be studied in
$\mathfrak C_2^{\rm APP}$, where \textup{(GMC-6.9)} permits exact computations.

\subsection{Explicit construction of quantum magic cubes and compatibility of their margins}

The paper on Latin squares proposes the development of quantum magic cubes.
The cyclic lemma admits an immediate extension.

\begin{definition}[Quantum Magic Cube]
A family
$C=(C_{ijk})_{0\le i,j,k<n}$ of effects in $M_s$ is a quantum magic cube
if every line parallel to any of the three axes sums to $I_s$.
\end{definition}

\begin{theorem}[Cyclic cube of a \POVM]
\label{p12-83-c22-s5:gmc:thm:cubo-ciclico}
For every \POVM $(E_t)_{t\in\mathbb Z/n\mathbb Z}$,
\[
 C_{ijk}=E_{i+j+k\bmod n}
 \tag{GMC-8.2}
\]
is a quantum magic cube of exterior size $n$ and interior size $s$.
\end{theorem}

\begin{proof}
Once two coordinates are fixed, as the third varies the index
$i+j+k$ traverses $\mathbb Z/n\mathbb Z$ once. Each line therefore sums to
$\sum_tE_t=I_s$. Positivity is inherited from \POVM.
\end{proof}

Applying the theorem to $E^{(9)}$ of \textup{(GMC-4.6)} and to $F$ of
\textup{(GMC-4.12)} yields explicit qutrit cubes of orders nine and twelve. The
former may assign the third axis to a phase coordinate; the latter uses the
complete nonadic atlas. To convert that axis into genuine HMT depth,
the layers must also satisfy
\[
 \rho_m(C^{m+1}_{ijk})=C^m_{ijk},
 \tag{GMC-8.3}
\]
and intertwine the TPK. The cyclic construction closes the operatorial cube;
genealogical fidelity remains an additional condition.

\subsection{Magic combinations of CP maps}

Another open direction of the theory consists in transferring the magic conditions
to completely positive maps through Choi--Jamiołkowski. If
$J(\Phi_{ij})\succeq0$ are the Choi matrices of a grid of CP maps,
the conditions
\[
 \sum_jJ(\Phi_{ij})=J(\id),
 \qquad
 \sum_iJ(\Phi_{ij})=J(\id)
 \tag{GMC-8.4}
\]
are a linear version of the magic normalizations. HMT provides channels of
reading, updating, and transport; the problem is to decompose them into a
grid that simultaneously preserves normalization, sheet, and route. The generalized
conic formalism can treat codomains that are not simply matrix
PSD cones.

\subsection{Universality criterion in the genealogical domain}

The framework of Gonda, Reinhart, Stengele, and De les Coves defines a simulator in
a \emph{target--context} category by means of
\[
 s:P\otimes C\longrightarrow T\otimes C,
 \tag{GMC-8.5}
\]
with a functional compiler $s_T:P\to T$ and a context reduction.
Universality is the existence of a lax reduction to the trivial simulator,
not the mere richness or recurrence of the system
\citep{gonda2024universality}.

An HMT instance may take:

\begin{itemize}
 \item $P$: prefixes, route programs, or seeds;
 \item $T$: mathematical and physico-metrological observables intended to be
realized;
 \item $C$: gauge, leaf, boundary conditions, and reading apparatus;
\item $s_T$: transformation of a genealogy into an output;
 \item $s_C$: transport of the context and residual memory.
\end{itemize}

Proving universality would require constructing those data and the lax reduction. A monotone of the contextual preorder suffices to refute a proposal. The framework's impossibility theorem states, in particular, that if the functional image of the compiler has a supremum strictly below the target space for some compatible monotone, the simulator is not universal.

\begin{bridgebox}
HMT suggests a stronger version: \emph{genealogically faithful universality}. In addition to reproducing the behavior of the target, the reduction must lift to a functor between categories of routes and preserve composition, return, and the enriched chronological register. Two behaviorally equivalent simulators may cease to be equivalent when preservation of provenance is required. This distinction offers a new and nontrivial instance for the categorical program.
\end{bridgebox}

The metrological unification demonstrated by HMT and categorical universality
are distinct claims. The former derives readers and relations within
the system; the latter requires the ability to simulate a declared class of targets and
contexts. Typing them separately strengthens both: it avoids requiring from metrology
a property that it does not need and turns universality into a falsifiable
theorem.

\subsection{Contextuality and Nonsemiclassicality}

The cyclic squares of \Cref{p12-83-c22-s3:gmc:chap:qms} are semiclassical. They therefore do not
resolve the open direction concerning contextuality or nonsemiclassical
QMS. APP, however, contains the noncommuting pair $P_\Sigma,P_\Pi$
and algebra \textup{(GMC-6.9)}. The correct route is to construct cell effects that
depend on both operators and prove that they admit no decomposition
\textup{(GMC-4.4)} with a \POVM on $S_n$.

The verification is a semidefinite problem: find
$Q_\sigma\succeq0$, $\sum_\sigma Q_\sigma=I$, and
\[
 A_{ij}=\sum_{\sigma:\,\sigma(i)=j}Q_\sigma.
 \tag{GMC-8.6}
\]
Feasibility certifies semiclassicality; infeasibility must be accompanied by
a dual witness. A noncommutative matrix or a visual graph does not replace
that certificate.


\section{Genealogical structure of APP--TRIT--TPK at the matrix level}
\label{p12-83-c22-s6:gmc:chap:hmt-profundidad}

This chapter isolates the structures involved in the transduction and keeps their types explicit. The order is deliberate: support, sum--product difference, ternary memory, category of routes, quantum realization, operatorial closure, and metrological reading. Only afterward will it be asked which of these structures form a magic square or an operator system.

\subsection{Common support of APP and dual evaluation via sum and product sheets}

Let
\[
 X_9=(\mathbb Z/9\mathbb Z)^2,
 \qquad
 \mathscr D=\{(1,0),(0,1),(-1,0),(0,-1)\}.
\]
The Cayley graph $\operatorname{Cay}(X_9,\mathscr D)$ has 81 vertices and four oriented cardinal edges from every vertex. The APP architecture contains all finite paths in that graph. The acronym is retained as a technical name and is not used here to fix a historical expansion. At length $r$ there are exactly $81\cdot4^r$ paths if returns and backtracking are allowed. There are not independent additive and multiplicative tables: there is a single support and two evaluations of the same point,
\[
 \Sigma(i,j)=\operatorname{dr}(i+j),
 \qquad
 \Pi(i,j)=\operatorname{dr}(ij),
 \qquad
 \operatorname{dr}(a)=1+((a-1)\bmod9).
 \tag{GMC-2.1}
\]
The table of $\Sigma$ is a classical Latin square of order nine, because
every translation permutes $\mathbb Z/9\mathbb Z$. The table of $\Pi$ is not:
only its rows and columns indexed by the six units
$\{1,2,4,5,7,8\}$ are permutations; $3,6,9$ fold classes.

\begin{resultbox}
APP already contains an exact Latin structure, but APP is not exhausted by it.
The sum leaf is Latin; the product leaf records the defect of Latinity of the
local ring; the paths preserve the consultation order, and the
residue--quotient lift preserves what modular reduction would conceal.
\end{resultbox}

For each cell, the residue and quotient are retained:
\[
 i+j=R^+_{ij}+9q^+_{ij},
 \qquad
 ij=R^\times_{ij}+9q^\times_{ij}.
 \tag{GMC-2.2}
\]
The certified totals are
\[
 \sum R^+=405,
 \quad \sum R^\times=459,
 \quad \sum q^+=45,
 \quad \sum q^\times=174,
\]
and therefore
\[
 2025-810=1215=(459-405)+9(174-45)=54+9\cdot129.
 \tag{GMC-2.3}
\]
The visible difference $54$ and the quotient difference $129$ are not two
names for the same datum. One resides in the published residues; the other, in
the memory required to recompose the integer value.

\subsection{Spectral hierarchy \(4\to2\to6\to54\) of the central APP block}

The central block of APP takes the form
\[
 B_c=\begin{pmatrix}7&2\\2&7\end{pmatrix}
     =7I+2R
     =9P_++5P_-,
 \qquad
 R=\begin{pmatrix}0&1\\1&0\end{pmatrix},
 \quad P_\pm=\frac{I\pm R}{2}.
 \tag{GMC-2.4}
\]
From this one obtains exactly the typed hierarchy induced by the
preceding spectral decomposition:
\[
 \underbrace{9-5}_{4}
 \ ;\ 
 \underbrace{2}_{\text{off-diagonal coupling}}
 \ ;\ 
 \underbrace{51-45}_{6}
 \ ;\ 
 \underbrace{9\cdot6}_{54}.
 \tag{GMC-2.5}
\]
The first figure is the primordial spectral gap; the second, the coupling
of the block; the third, the aggregate excess of the nonunitary sector; the fourth,
its transport to the global scale of the table. They must not be collapsed on account of their
arithmetic kinship.

Since $\det B_c=45$, the normalization has determinant one and belongs to the
connected hyperbolic group:
\[
 \frac{B_c}{\sqrt{45}}
 =\exp(\eta R),
 \qquad
 \eta=\frac12\log\frac95,
 \qquad
 \frac59=e^{-2\eta}.
 \tag{GMC-2.6}
\]
The same spectral pair thus admits a contractive reading and a
Lorentzian reading. This is the precise connection with convexity and the
hyperbolic interior that the bridge must preserve: the contraction $5/9$ is the shadow
of the separation of the two idempotents $P_\pm$, not a verbal analogy with
curvature.

\subsection{TRIT: visible state, orientation, carry, and genealogical memory}

TRIT is the oriented signature $\{-1,0,+1\}$ obtained from $\mathbb F_3$ through
$0\mapsto0$, $1\mapsto+1$, $2\mapsto-1$. To type its local lift,
we define the lifted alphabet
\[
 \mathcal L_1=\{\nu=r+3c:r,c\in\{-1,0,+1\}\}
 =\{-4,-3,\ldots,4\},
 \qquad
 \widehat\tau(\nu)=(r,c).
 \tag{GMC-2.7}
\]
The decomposition $\nu=r+3c$ is unique in $\mathcal L_1$; $r$ is the visible
residue and $c$ the carry. The symbol $\nu$ is a lifted amplitude and is not
identified, in the absence of further data, with the sum of two visible trits. In particular,
the amplitudes $3,0,-3$ publish the same zero residue and yield three distinct
antecedents,
\[
 0^+=(0,+1),\qquad 0^0=(0,0),\qquad0^-=(0,-1).
 \tag{GMC-2.8}
\]
The complete state may contain
\[
 x=(t,\mu,o,h,c,\sigma,\lambda,g),
 \tag{GMC-2.9}
\]
where $t$ is the visible signature, $\mu$ the leaf, $o$ the orientation, $h$ the
extension, $c$ the carry, $\sigma$ the boundary, $\lambda$ the route, and $g$ the
nonadic phase. The projection $x\mapsto t$ eliminates all the other
coordinates. This is the structural point that HMT contributes to the theory of
dilations: two equal observable effects may possess distinct
genealogies.

The three signatures admit the real algebras
\[
 \mathbb A_\tau=\mathbb R[J]/(J^2+\bar\tau),
\]
of elliptic type ($J^2=-I$), parabolic type ($J^2=0$), and hyperbolic type
($J^2=+I$). In the last case, $P_\pm=(I\pm J)/2$ separate two eigendirections.
The TRIT is not merely the sign: it is the minimal observable of a transport
that preserves the label, orientation, and memory.

\subsection{TPK: category of memoryful paths}

TPK, \emph{Topological Phase Kernel}, resides on the free path category of APP and on its bivariate sum--product reading. Two paths that end in the same cell may differ in cardinal word, winding, sheet, quotient, and holonomy. TPK preserves that difference and composes routes associatively. The carry cocycle
\[
 \kappa_9(a,b)=
 \frac{\operatorname{dr}(a)+\operatorname{dr}(b)
       -\operatorname{dr}(a+b)}9
\]
satisfies
\[
 \kappa_9(a,b)+\kappa_9(a+b,c)
 =\kappa_9(b,c)+\kappa_9(a,b+c),
 \tag{GMC-2.10}
\]
and guarantees that the recovery does not depend on how a concatenation is parenthesized.

The quotient phase--length has objects $(b,r)\in\mathbb Z/12\mathbb Z
\times\{0,\ldots,8\}$ and arrows
\[
 (b,r)\xrightarrow{\ell}
 \left(b+\left\lfloor\frac{r+\ell}{9}\right\rfloor,
       r+\ell\bmod9\right).
 \tag{GMC-2.11}
\]
After nine unit iterations ($\ell=1$), $r$ returns and one
dodecaphasic phase advances; after $108=9\cdot12$ unit iterations, the quotient returns.
The enriched state of TPK may not
have returned, because it preserves additional memory. The minimal predictive
quotient of the published observables has 36 states and fibers of cardinality
three. This fact makes visible the difference between equality of output and
equivalence of history.

\begin{intuitionbox}
A magic square simultaneously organizes rows and columns. TPK simultaneously organizes position, sheet, phase, and composition. The correct translation does not consist in calling every HMT table “quantum,” but in constructing a positive reader that preserves the APP identities and declaring, coordinate by coordinate, which memory it preserves and which memory it compresses.
\end{intuitionbox}

\subsection{Hilbertian realization and quantum structure}

The finite quantum closure of HMT constructs, for each depth
$d$, the space
\[
 \mathcal H_d=\ell^2((\mathbb Z/9\mathbb Z)^d),
 \qquad
 \mathcal A_d=B(\mathcal H_d)\cong M_{9^d}(\mathbb C).
 \tag{GMC-2.12}
\]
For $u,v\in(\mathbb Z/9\mathbb Z)^d$, the nonadic translations and phases
act by means of
\[
 X_u|x\rangle=|x+u\rangle,
 \qquad
 Z_v|x\rangle=\omega^{v\cdot x}|x\rangle,
 \qquad \omega=e^{2\pi i/9},
\]
and satisfy
\[
 Z_vX_u=\omega^{u\cdot v}X_uZ_v.
\]
The complete family---equivalently, the $d$ coordinate pairs
$X_j,Z_j$---generates $\mathcal A_d$ and acts irreducibly. For the primitive
central character fixed by $\omega$, the finite Heisenberg representation
is determined up to unitary equivalence. In this realization, states,
self-adjoint observables, spectral measures, Born probabilities, unitary
dynamics, and Kraus channels are constructed; the composition of registers
is transported monoidally to the tensor product. The finite Hilbert-space
representation is not presupposed: it is constructed explicitly and transports
the genealogy of the system.

In the qutrit quotient there appear four maximal abelian subalgebras of the
three-dimensional Weyl system. Their four spectral bases are
mutually unbiased:
\[
 \Tr(P_{a,k}P_{b,\ell})=
 \begin{cases}
  \delta_{k\ell},&a=b,\\[2pt]
  1/3,&a\ne b,
 \end{cases}
 \qquad a,b\in\mathbb P^1(\mathbb F_3).
 \tag{GMC-2.13}
\]
The twelve minimal projectors $P_{a,k}$ are the operatorial atlas that is later
calibrated with the twelve directions of
$\mathbb P^1(\mathbb Z/9\mathbb Z)$.

\subsection{From the nonadic tower to the operatorial continuum}

The unital inclusion, the marked branch, the UHF limit, and the closures of
types \(I_\infty\) and \(II_1\) are constructed in
\cref{chap:integracion-vn64-operatorial}. Their formulas are premises of the matrix bridge and determine
the bigrading.

The operatorial tower prolongs the finite charts through the
unital inclusions preserving the normalized trace
\[
 \iota_d:M_{9^d}(\mathbb C)\longrightarrow M_{9^{d+1}}(\mathbb C),
 \qquad a\longmapsto a\otimes I_9.
 \tag{GMC-2.14}
\]
Indeed, if $\tau_n=\Tr/n$, then
$\tau_{9^{d+1}}\circ\iota_d=\tau_{9^d}$; the ordinary trace, by contrast, is
multiplied by nine. 
Its inductive limit is the UHF algebra of type $9^\infty=3^\infty$. The closure
in the GNS representation of the unique trace is the hyperfinite factor
$\Rhyper$ of type $\mathrm{II}_1$. The marked inclusions
$a\mapsto a\otimes p_j$ follow one branch and lead to the corridor of compact operators,
of type $\mathrm I_\infty$. The result therefore distinguishes preserving
all branches simultaneously from selecting a particular history.

In the diagonal corridor, the finite sets of prefixes form an inverse system
$X_{m+1}\to X_m$. The space of infinite histories is
$X_\infty=\varprojlim X_m$, while its function algebras form the
opposite inductive system. This duality of directions ---inverse for
histories, direct for observables--- will be the exact support of the new
bigraduation.

\subsection{Double projection and metrology}

HMT does not identify the published real number with the state that generates it. A
compatible family of cylinders produces an Archimedean reading; another
projection organizes the exceptional output and its symmetries. Numerical
reality appears as the shadow of a double projection of a richer state.
The monodromic generation theorems for $\pi$, $e$, $\varphi$, and $\alpha$,
the action and
gravitation readers, the mass law, the mixing and CP-violation sectors, and the
exceptional genealogy apply in their declared domains. The bridge makes
explicit the operator objects that transport their conclusions.

The Planck identity uses the same block $B_c$, its
transport in $SO^+(1,1)$, and the self-dual identity of the mass line. The
decisive mathematical fact for this translation is that ‘physical observable’ and
‘generating history’ are different types. A UCP map can publish
the former; the coordinate $m$ and the TPK category preserve the latter.

\begin{sourcebox}
The premises of the bridge are the unital inclusion
$a\mapsto a\otimes I_9$, the block $B_c=9P_++5P_-$, the self-dual identity of
the mass line, and periodicity of order \(108\). Each is used with
the domain and hypotheses of its theorem; the matrix conclusions are
deduced by means of the maps written in this chapter.
\end{sourcebox}


% El capítulo se distribuye en dos dependencias causales. La unión de los dos
% selectores reproduce todas sus líneas científicas exactamente una vez.
\ifdefined\HMTGMCContinuumMPS
\section{Genealogical Limits, Transfer Channels and Stinespring Memory}
\label{p12-83-c22-s7:gmc:chap:continuo-mps}

The HMT continuum and the continuum limit of matrix product states (MPS)
answer different questions. The former prolongs finite histories and
then reads them; the latter refines a tensor network and requires its transfer
channel to admit compatible roots. The comparison is fruitful only if
the channel connecting them is constructed.

\subsection{Typed separation of the genealogical limit and the MPS limit}

In the HMT corridor, an infinite history is an element of
$X_\infty=\varprojlim X_m$. Under the finiteness, surjectivity, nested-compactum,
contraction, and coherent-covering conditions proved in the continuum
appendix, the quotient by the Archimedean reading reconstructs the declared
real compactum. The complete state and its real value remain distinct
objects.

For a translationally invariant MPS, a tensor
$A=\{A^i\in M_D\}_{i=1}^d$ determines the transfer channel
\[
 E_A=\sum_i A^i\otimes\overline{A^i}.
 \tag{GMC-7.1}
\]
De las Cuevas, Schuch, Pérez-García, and Cirac prove that an MPS in
irreducible form has a continuum limit if and only if $E_A$ is infinitely
divisible; equivalently, there exist a projector channel $P$ and a Lindblad
generator $L$ such that
\[
 E_A=Pe^L,
 \qquad P^2=P,
 \qquad PLP=PL.
 \tag{GMC-7.2}
\]
It has a coarse continuous limit if some power $E_A^r$ is infinitely
divisible \citep{delascuevas2018continuum}.

\begin{controlbox}
The existence of the HMT inverse limit does not imply \textup{(GMC-7.2)}. To apply the MPS theorem, one must exhibit a uniform local tensor, identify its transfer channel, and verify infinite divisibility. Conversely, \textup{(GMC-7.2)} does not by itself preserve sheet, carry, route, or monodromy.
\end{controlbox}
\fi

\ifdefined\HMTGMCOperatorialTower
\subsection{Nonadic refinement through conditional expectation}
\label{p12-83-c22-s7:gmc:sec:refinamiento-condicional}

In the tower
$\mathcal A_{m+1}=\mathcal A_m\otimes M_9$, the conditional expectation
\[
 \mathbb E_m=\id\otimes\tau_9:
 \mathcal A_{m+1}\longrightarrow\mathcal A_m,
 \qquad \tau_9(b)=\frac1{9}\Tr(b),
 \tag{GMC-7.3}
\]
is UCP and satisfies
\[
 \mathbb E_m(a\otimes I_9)=a.
 \tag{GMC-7.4}
\]
On the diagonal, $D_{m+1}=C(X_m\times\mathbb Z_9)$ reduces to the
uniform average over the last digit. The branch evaluation
\[
 \kappa_{m,a}(f)(x)=f(x,a)
 \tag{GMC-7.5}
\]
is another retraction, this time multiplicative on the diagonal. Averaging all
branches and selecting one branch are distinct operations, parallel to the
inclusions $a\otimes I_9$ and $a\otimes p_a$ of the operatorial closure.

The expectation \textup{(GMC-7.3)} has a dilation with an explicit record. Let
\[
 W_a\psi=\frac13\,\psi\otimes e_a,
 \qquad a=0,\ldots,8.
 \]
Then
\[
 \mathbb E_m(X)=\sum_{a=0}^8W_a^*XW_a.
 \tag{GMC-7.6}
\]
In Stinespring form,
\[
 \pi(X)=X\otimes I_9,
 \qquad
 V\psi=\frac13\sum_{a=0}^8(\psi\otimes e_a)\otimes e_a,
 \qquad
 \mathbb E_m(X)=V^*\pi(X)V.
 \tag{GMC-7.7}
\]
The final auxiliary system (\emph{ancilla}) labels exactly the digit discarded by averaging.

\begin{bridgebox}
The formula \textup{(GMC-7.7)} is a canonical nonadic memory in the sense of
Stinespring: what is lost when applying $\mathbb E_m$ remains orthogonally
registered. It is not yet identified with TPK memory. That identification
requires an intertwiner that also transports route, leaf, and dynamics.
\end{bridgebox}

\subsection{Stinespring's Memoir and Non-Uniform Growth}
\label{p12-83-c22-s7:gmc:sec:memoria-no-uniforme}

The \Cref{p12-83-c22-s4:gmc:thm:dilatacion-uhf} provides a common, generally
infinite-dimensional representation for every compatible family. A fixed
finite ancillary system that purifies all levels cannot in general be required. Indeed, the
compatible family of normalized traces on $M_{9^m}$ has rank $9^m$ at
level $m$, and any purification of the corresponding state requires
an ancillary dimension of at least $9^m$. Therefore,
\[
 \text{UCP compatibility at every depth}
 \quad\centernot\Longrightarrow\quad
 \text{uniform finite auxiliary memory}.
 \tag{GMC-7.8}
\]
Memory growth is not a defect of the construction: it quantifies the
amount of genealogy that one requires to be preserved.

The dynamic fidelity of this memory and its comparison with the limit of
matrix product states are developed in
\cref{p12-83-c22-s7:gmc:sec:fidelidad-dinamica}, without repeating the UHF tower or the Stinespring dilation
demonstrated here.
\fi

\ifdefined\HMTGMCContinuumMPS
\subsection{Dynamic fidelity equation}
\label{p12-83-c22-s7:gmc:sec:fidelidad-dinamica}

The trace-preserving expectation \(\id\otimes\tau_9\), its Stinespring
dilation, and the general impossibility of a uniform finite auxiliary
memory have been constructed in
\cref{p12-83-c22-s7:gmc:sec:refinamiento-condicional,p12-83-c22-s7:gmc:sec:memoria-no-uniforme}. On
that operatorial premise, and without abbreviating its proofs, the
additional condition of genealogical fidelity is now formulated.

Let $T_m^{\rm TPK}$ be the transport between records at depth $m$, and let
$J_m$ be the encoding of those records in the dilation space. The
Stinespring memory is TPK-faithful if there exists a transport
$\widetilde T_m$ in the ambient space such that
\[
 J_{m+1}T_m^{\rm TPK}=\widetilde T_mJ_m
 \tag{GMC-7.9}
\]
and the observational compressions satisfy the corresponding naturality.
If two distinct histories are identified under $J_m$, the dilation may be
operatorially valid but not genealogically faithful.

The equation \textup{(GMC-7.9)} is falsifiable: it suffices to have a pair of routes that TPK
distinguishes and the environment identifies, or a composition whose image does not coincide
with the composition of the images.

\subsection{Compatibility Criterion for Product State Limits}

To apply \textup{(GMC-7.2)} to the HMT continuum, five steps must be completed:

\begin{enumerate}
 \item choose a uniform local cell and construct matrices
 $A^i\in M_D$;
 \item prove that nonadic cell blocking intertwines the evolution
 with a refinement isometry;
 \item compute the channel $E_A$ and its irreducible form;
 \item resolve whether $E_A$ or some power $E_A^r$ is infinitely
divisible;
 \item transport the leaf, carry cocycle, and
 monodromy to the limit, rather than preserving only the tensor state.
\end{enumerate}

If the channel were unitary, $E_A=\Ad_U$, it would have unitary roots of every
order by functional calculus. But it would still be necessary to prove that this channel is
the transfer channel associated with an HMT tensor and that its roots respect the
genealogical refinement. If a projective part $P$ appeared, HMT memory would offer
a concrete interpretation of the sectors that survive scale grouping.

\subsection{Tensorial decompositions with symmetry}

The work of De las Cuevas with Netzer and collaborators on
approximate tensor decompositions and on simplicial complexes
provides another point of entry \citep{delascuevas2021approximate,
delascuevas2024simplicial}. APP supplies a finite complex with an action of
translations; TPK supplies labels and compatibilities along routes. An
HMT tensor decomposition should state:

\begin{itemize}
 \item which tensor encodes a cell or a transition;
 \item which group preserves the labels and which part acts only after
 reduction;
 \item whether the factorization is exact, approximate, or only asymptotic;
 \item which rank remains stable as $m$ increases.
\end{itemize}

The bidirectional opportunity is concrete. Its formalism provides criteria of
factorization, invariance, and limit; HMT provides a refinement system with
fibers, carry, and monodromy that can produce examples in which a visible
factorization exists but no genealogically faithful factorization
preserves the route.

\begin{workbox}
A subsequent paper may study "infinite divisibility with root
memory": not only whether $E=E_k^k$, but whether a coherent tower of
roots $E_{k\ell}^\ell=E_k$ equipped with TPK records can be chosen. This condition is
stronger than divisibility of the bare channel and constitutes a natural
extension of the MPS problem from HMT.
\end{workbox}
\fi


\section{Completely positive matrix realization of the nonadic fibration of observables}
\label{p12-83-c22-s8:gmc:chap:simplectica}

Symplectic geometry is not part of the published theorems on
De las Cuevas--Netzer magic squares. The legitimate link passes through a
neighboring, standard route: Weyl operators, discrete phase space and
mutually unbiased bases \citep{gibbons2004phase,planat2007rings}. In HMT,
that route acquires a nonadic lift and fiber memory.

\subsection{The alternating form and the Weyl commutation}

On $V_3=\mathbb F_3^2$ we define
\[
 \langle(q,p),(q',p')\rangle_{\rm sp}=qp'-pq'.
 \tag{GMC-5.1}
\]
If $X|k\rangle=|k+1\rangle$,
$Z|k\rangle=\omega^k|k\rangle$, and $\omega=e^{2\pi i/3}$, the operators
$W_{q,p}=X^qZ^p$ satisfy, under these conventions,
\[
 W_uW_v=\omega^{-\langle u,v\rangle_{\rm sp}}W_vW_u.
 \tag{GMC-5.2}
\]
Therefore, two operators commute exactly when their labels are
orthogonal for the alternating form. In dimension two, the maximal
isotropic sublines are the four points of
\[
 \mathbb P^1(\mathbb F_3)=\{[1:0],[0:1],[1:1],[1:2]\}.
 \tag{GMC-5.3}
\]
Each generates a maximal abelian context of $M_3(\mathbb C)$; its
spectral projectors constitute a \PVM. Contexts associated with distinct
lines produce the four MUB bases of \textup{(GMC-2.13)}.

\begin{proposition}[Qutrit decomposition by directions]
\label{p12-83-c22-s8:gmc:prop:descomposicion-masa-qutrit}
Let $\mathcal D_a$ be the four diagonal subalgebras determined by the
lines of \textup{(GMC-5.3)} and $\mathcal D_a^0$ their traceless parts. Then
\[
 M_3(\mathbb C)
 =\mathbb CI\oplus
   \bigoplus_{a\in\mathbb P^1(\mathbb F_3)}\mathcal D_a^0
 \tag{GMC-5.4}
\]
as an orthogonal sum for the Hilbert--Schmidt product. The dimension
count is $9=1+4\cdot2$.
\end{proposition}

\begin{proof}
Each $\mathcal D_a^0$ has dimension two. Unbiasedness implies
$\Tr(A^*B)=0$ for $A\in\mathcal D_a^0$ and
$B\in\mathcal D_b^0$, $a\ne b$. All are orthogonal to $I$. The
dimensions sum to nine, the dimension of $M_3(\mathbb C)$.
\end{proof}

This identity explains why the twelve projections are not twelve arbitrary
independent pieces: they are four resolutions of the identity whose centered
parts decompose the entire qutrit operator space.

\subsection{The lift \texorpdfstring{$12\to4\times3$}{12 to 4 by 3}}

On the modulo nonadic $V_9=(\mathbb Z/9\mathbb Z)^2$ it preserves the forma
alternante
\[
 \langle(a,b),(c,d)\rangle_9=ad-bc\pmod9.
 \tag{GMC-5.5}
\]
The projective line $\mathbb P^1(\mathbb Z/9\mathbb Z)$ consists of
the primitive vectors modulo multiplication by a unit. It can
be represented as
\[
 \{[1:t]:t\in\mathbb Z/9\mathbb Z\}
 \ \sqcup\ 
 \{[3u:1]:u\in\mathbb Z/3\mathbb Z\},
 \tag{GMC-5.6}
\]
and has $9+3=12$ points. Reduction modulo three gives the diagram
\[
 \mathbb P^1(\mathbb Z/9\mathbb Z)
 \xrightarrow{\ r\ }
 \mathbb P^1(\mathbb F_3),
 \qquad |r^{-1}(a)|=3.
 \tag{GMC-5.7}
\]

The basis $a$ of a qutrit context and the outcome $k$
of its \PVM likewise form a set $4\times3$. A
\emph{calibration} is a bijection
\[
 \Xi:\mathbb P^1(\mathbb Z/9\mathbb Z)
 \longrightarrow
 \{P_{a,k}:a\in\mathbb P^1(\mathbb F_3),\ k\in\mathbb F_3\}
 \tag{GMC-5.8}
\]
that sends each fiber of $r$ to the three outcomes of the context lying over
$a$.

\begin{controlbox}
The fibration is canonical; assigning the internal order of each fiber to the
three outputs $k$ is a gauge. A bijection is not yet an
equivariance. To promote $\Xi$ to an intertwiner, one must prove that the
relevant linear, projective, and antiunitary actions are transported in
every fiber. The preliminary gauge fails an antiunitary check in one
fiber; that failure is preserved as a falsifier, not concealed.
\end{controlbox}

\subsection{Depth, carry, and memory over the discrete phase space}

The Weyl geometry over $\mathbb F_3^2$ explains commutation, contexts, and
MUB. HMT adds three levels that are not recovered from the symplectic form by
itself:

\begin{enumerate}
 \item the lift of a ternary direction to three points of
 $\mathbb P^1(\mathbb Z/9\mathbb Z)$;
 \item the sheet, carry and orientation that distinguish states with the same reduction;
 \item the TPK route that transports those coordinates and can return to the
 phase without returning to the state.
\end{enumerate}

The potential contribution to the qutrit literature is not “another construction of
four MUBs,” which is already known. It is a geometry of \emph{context roots}: each
ternary context has three nonadic antecedents, and those antecedents
may carry dynamics and monodromy. The order-twelve square of
\Cref{p12-83-c22-s3:gmc:thm:qms-12-3} is the first normalized shadow of that atlas.

\subsection{Typed separation between the qutrit overlap
\texorpdfstring{$1/3$}{1/3} and the HMT reader
\texorpdfstring{$1/\pi$}{1/pi}}

In a pair of MUB qutrit bases,
\[
 \Tr(P_{a,k}P_{b,\ell})=\frac13.
 \tag{GMC-5.9}
\]
In the HMT reader defined here, the relation between a real sheet and a
volumetric reading can take the quotient
\[
 \frac{\pi}{\pi^2}=\frac1\pi.
 \tag{GMC-5.10}
\]
The two expressions are not interchangeable: $1/3$ is an overlap between
projectors in dimension three; $1/\pi$ is a scalar weight of another reader.
The difference does not prevent their coordination. For a fixed \PVM
$\{P_0,P_1,P_2\}$, we define a measurement with nine effects:
\[
 F_k^{(\pi)}=\frac3\pi P_k\quad(k=0,1,2),
 \qquad
 F_r^{(\pi)}=\frac{1-3/\pi}{6}I_3\quad(r=3,\ldots,8).
 \tag{GMC-5.11}
\]
It is a \POVM because $3/\pi<1$ and
\[
 \sum_{r=0}^{8}F_r^{(\pi)}
 =\frac3\pi I_3+\left(1-\frac3\pi\right)I_3=I_3.
 \tag{GMC-5.12}
\]
If $Q$ is a rank-one projector of a distinct MUB context,
\[
 \Tr(QF_k^{(\pi)})
 =\frac3\pi\Tr(QP_k)
 =\frac3\pi\cdot\frac13
 =\frac1\pi.
 \tag{GMC-5.13}
\]

\begin{bridgebox}
The formula \textup{(GMC-5.11)} exactly transforms the MUB probability $1/3$ into the reader $1/\pi$ through the gauge $3/\pi$, without identifying the two numbers. The six residual effects preserve the unit and complete a \POVM of nine outcomes, from which another cyclic \QMS $(9,3)$ is obtained. This construction does not yet identify the geometric provenance of the residual mass; it fixes the exact operator type that it must have.
\end{bridgebox}

\subsection{Symplectic group and equivariance equation}

In dimension two over $\mathbb F_3$, the symplectic group coincides with
$SL(2,\mathbb F_3)$ and permutes the four lines. Its action lifts
projectively, by means of the Clifford group, on the \PVM qutrits. Let
$G_{\rm HMT}$ be a group or groupoid of transformations of the nonadic atlas and
$U_g$ the candidate unitary or antiunitary operator in the qutrit chart. The
condition that turns a calibration into a dynamical bridge is
\[
 \Xi(g\cdot\ell)=U_g\,\Xi(\ell)\,U_g^*,
 \qquad
 g\in G_{\rm HMT},\quad
 \ell\in\mathbb P^1(\mathbb Z/9\mathbb Z).
\tag{GMC-5.14}
\]
If $g$ is an arrow and not merely a global symmetry, the equation must include
the source and target and is formulated as a natural transformation. The verification of
\textup{(GMC-5.14)} is the step that separates a fortunate labeling from a
structural transduction.

\begin{sourcebox}
The symplectic and MUB identities rely on the finite-phase literature
\citep{gibbons2004phase,planat2007rings}. They are not attributed to the papers on
magic squares by De las Cuevas--Netzer. The nonadic fibration, calibration,
and route readers constitute the original contribution of the
HMT construction studied here.
\end{sourcebox}
